\documentclass[10pt,a4paper]{amsart}
\usepackage[margin=19mm]{geometry}
\usepackage{amsmath,amssymb,mathtools}
\usepackage[scr=rsfs,cal=euler]{mathalfa}
\usepackage{xfrac}
\usepackage[unicode,hidelinks,bookmarks=false]{hyperref}
\newcommand{\ie}{i.e.\ }
\newcommand{\cf}{cf.\ }
\newcommand{\resp}{respectively\ }
\newcommand{\Subsection}{\subsection}
\providecommand{\nobreakdash}{\nobreak\hbox{-}}
\setlength{\emergencystretch}{2em}
\multlinegap0pt
\usepackage{amssymb}
\usepackage{float}
\newtheorem{lemma}{Lemma}
\newcommand{\CF}{\mathcal{F}}
\newcommand{\CN}{\mathcal{N}}
\newcommand{\CO}{\mathcal{O}}
\newcommand{\X}{\times}
\newcommand{\fl}{\mathfrak{l}}
\newcommand{\g}{\mathfrak{g}}
\title{Character Sheaves on Reductive Lie Algebras in Positive Characteristic}
\author{Tong Zhou}
\date{Source: arXiv:2404.19210v2 (CC BY 4.0)}
\begin{document}
\maketitle

\noindent\emph{Source and licence.} Character Sheaves on Reductive Lie Algebras in Positive Characteristic. Version arXiv:2404.19210v2 (CC BY 4.0), distributed by arXiv under the Creative Commons Attribution 4.0 International licence, http://creativecommons.org/licenses/by/4.0/, as recorded in the arXiv licence metadata for this exact version and shown on the abstract page https://arxiv.org/abs/2404.19210v2. Full licence notice: \texttt{licenses/arxiv-2404.19210-CC-BY-4.0.txt}.\par\smallskip

\noindent\textit{Editorial scope.} Counted unit: the article's lemma lem\_form\_of\_SS (source0.tex lines 1477--1480). The article's complete original proof of it is delivered with it (source0.tex lines 1482--1488, one \texttt{proof} environment of 7 lines). No mathematics is written, repaired or re-derived: the statement and the proof are the article's own lines, in the article's own notation, and no retained line is substituted or re-flowed.  Retained uncounted as the source-local context the proof needs: lines 1494--1496.  Nothing was omitted from any retained span, so the package carries no omission record.  No retained line contains a citation, so the wrapper carries no bibliography and no bibliography entry.  No retained line contains a figure environment, an image-inclusion command or drawing code, so no image is delivered and the package declares no asset dependency.  Editorial formatting only, in addition: an AMS \texttt{amsart} preamble that restates the article's own declarations and macros used by the retained lines, namely the environments \texttt{lemma} on separate counters, together with the standard packages those lines need.  The retained environments print in this document as Lemma 1 in order of appearance, whereas the article prints the same items with the numbering of its own full document.  The complete native source of the article as distributed in the arXiv e-print for this exact version is bundled unchanged as \texttt{sources/158369/source0.tex}.
\begin{lemma}\label{lem_conormal_comp}
    Let $x\in\g$, $L=Z_G(x_s)^\circ$. Then, $L$ is a Levi subgroup and $x\in Z_r(\fl)+\CO$ for some nilpotent orbit $\CO$ of $L$. Let $S_{(L,\CO)}$ be the corresponding Lusztig stratum. Then $(T^*_{Z_r(\fl)+\CO}\fl)_x=Z_{[\fl,\fl]}(x_n)$,  $(T^*_{S_{(L,\CO)}}\g)_x=Z_{[\fl,\fl]}(x_n)$.
\end{lemma}

\begin{lemma}\label{lem_form_of_SS}
    1) $\Lambda$ is closed, Lagrangian, and is contained in $\cup_{all} (T^*_{S_{(L,\CO)}}\g)$, where the union ranges exactly once through each $G$-conjugacy class of $(L,\CO)$ pairs. It follows, for dimensional reasons, that $\Lambda$ is of the form $\cup_{some} (\overline{T^*_{S_{(L,\CO)}}\g})$, for some classes of $(L,\CO)$.\\
    2) If $\CF\in Perv_G(\g)$ has nilpotent singular support, then $SS(\CF)\subseteq \Lambda$. It follows, for dimensional reasons, that $SS(\CF)$ is of the form $\cup_{some} (\overline{T^*_{S_{(L,\CO)}}\g})$, for some classes of $(L,\CO)$.
\end{lemma}

\begin{proof}
    1) $\Lambda$ is closed because it is defined by a closed condition. To see it is Lagrangian, first note that, in general, for any $x\in\g$, $(T^*_{G(x)}\g)_x=Z_\g(x)$.\footnote{Proof: $(T^*_{G(x)}\g)_x=((TG(x))_x)^\perp=[\g,x]^\perp=Z_\g(x)$, where in the second equality we have used the fact that, under our assumption on $p$, $Lie(Z_G(x))=Z_\g(x)$ (see Conventions).} As $\Lambda$ coincides with $=\{(x,y)\in\g\X\CN_G\,|\, x\in Z_\g(y)\}\subseteq \g\X\g$. View $\g\X\g$ as the cotangent bundle of the second $\g$ factor, $\Lambda$ is then precisely $\cup_{\CO\subseteq\CN_G}\g$. This shows it is Lagrangian.\\
    
    We now show $\Lambda\subseteq\cup_{all} (T^*_{S_{(L,\CO)}}\g)$. For any $x\in\g$, let $L=Z_G(x_s)^\circ$. Then, by Lemma \ref{lem_conormal_comp}, $x\in Z_r(\fl)+\CO$ for some nilpotent orbit $\CO$ of $L$, and $(T^*_{S_{(L,\CO)}}\g)_x=Z_{[\fl,\fl]}(x_n)$. As $Z_\g(x)\cap\CN_G=Z_\fl(x)\cap\CN_G=Z_\fl(x)\cap\CN_L=Z_{[\fl,\fl]}(x_n)\cap\CN_L$ (we used $Z_\g(x)=Z_\fl(x)$, see the proof of Lemma \ref{lem_conormal_comp}), the claim follows.\\

    2) As $\CF$ is $G$-equivariant, $SS(\CF)\subseteq \cup_{x\in \g}((T^*_{G(x)}\g)_x)=\cup_{x\in \g}(Z_\g(x)_x)$. As $SS(\CF)\subseteq \g\X\CN_G$, $SS(\CF)\subseteq\{(x,y)\in\g\X\CN_G\,|\, y\in Z_\g(x)\}=\Lambda$.
\end{proof}

\end{document}
